\uuid{dX9G}
\titre{Résolution d'une équation différentielle dans $\mathcal{D}'_+$}

\niveau{M1}
\module{Analyse}
\chapitre{Distributions}
\sousChapitre{Équations différentielles causales}
\theme{Second membre composé de dérivées de Dirac}
\auteur{Q. Liard}
\datecreate{ 2026-09-03}
\organisation{CentraleSupélec}
\difficulte{4}

\contenu{

\texte{
Déterminer la distribution $T\in\mathcal{D}'_+$ qui vérifie
$$
3T''+5T'+2T=3\delta'+\delta''.
$$
}


\question{Déterminer explicitement le signal $T$.}

\reponse{
Notons
$$
D=3\frac{d^2}{dt^2}+5\frac{d}{dt}+2.
$$
On commence par chercher une solution fondamentale causale $G$ de
$$
DG=\delta.
$$
Résolvons le problème de Cauchy
$$
\begin{cases}
3z''+5z'+2z=0,\\
z(0)=0,\\
z'(0)=1.
\end{cases}
$$
L'équation caractéristique est
$$
3r^2+5r+2=(3r+2)(r+1)=0.
$$
Ses racines sont $-2/3$ et $-1$. Ainsi,
$$
z(t)=Ae^{-t}+Be^{-\frac{2t}{3}}.
$$
Les conditions initiales donnent
$$
A+B=0,
\qquad
-A-\frac23B=1.
$$
On obtient $A=-3$ et $B=3$, donc
$$
z(t)=3\left(e^{-\frac{2t}{3}}-e^{-t}\right).
$$
Comme le coefficient dominant de $D$ vaut $3$, une solution fondamentale est
$$
G(t)=\frac13H(t)z(t)
=H(t)\left(e^{-\frac{2t}{3}}-e^{-t}\right).
$$
En effet, $DG=\delta$.

La solution cherchée est alors
$$
T=G\ast(3\delta'+\delta'')=3G'+G''.
$$
Posons
$$
q(t)=e^{-\frac{2t}{3}}-e^{-t},
\qquad G=Hq.
$$
On a $q(0)=0$ et
$$
q'(t)=e^{-t}-\frac23e^{-\frac{2t}{3}},
\qquad
q'(0)=\frac13.
$$
Donc
$$
G'=H\left(e^{-t}-\frac23e^{-\frac{2t}{3}}\right)
$$
et
$$
G''=\frac13\delta
+H\left(\frac49e^{-\frac{2t}{3}}-e^{-t}\right).
$$
Par conséquent,
$$
T
=3G'+G''
=
\frac13\delta
+H\left(
3e^{-t}-2e^{-\frac{2t}{3}}
+\frac49e^{-\frac{2t}{3}}-e^{-t}
\right).
$$
Après simplification,
$$
\boxed{
T
=
\frac13\delta
+H(t)\left(
2e^{-t}-\frac{14}{9}e^{-\frac{2t}{3}}
\right)
}.
$$
}



}
